\subsection{克利福德代数的基.}

命题 3. --- 设 Q 与 Q' 是两个二次型，F 是 E 上一个双线性型，使得对一切 x ∈ E 有 Q'(x) = Q(x) + F(x, x)。映射 λ\_F 把理想 I(Q') 映到理想 I(Q) 上，并定义 A-模 C(Q') 到 A-模 C(Q) 上的一个同构（记作 \(\overline{\lambda_{\mathbf{F}}}\)）。

  由于 \(\lambda_{\mathbb{F}}\) 是一个双射，其逆双射是 \(\lambda_{-\mathbb{F}}\)（引理 3），只需证明包含关系 \(\lambda_{\mathbb{F}}(\mathrm{I}(\mathrm{Q}^{\prime})) \subset \mathrm{I}(\mathrm{Q})\)。由于 \(\mathrm{I}(\mathrm{Q})\) 是在 \(i_x^{\mathbb{F}}\) 下稳定的左理想（引理 1），(8) 表明由 \(u \in \mathrm{T}(\mathrm{E})\)（使 \(\lambda_{\mathbb{F}}(u) \in \mathrm{I}(\mathrm{Q})\)）组成的集是一个左理想。因此只需证明：对任意 \(u \in \mathrm{T}(\mathrm{E})\) 与 \(x \in \mathrm{E}\)，有 \(\lambda_{\mathbb{F}}(x \otimes x \otimes u - \mathrm{Q}'(x)u) \in \mathrm{I}(\mathrm{Q})\)。现在，由 (8) 与引理 1，我们有

\[
\lambda_ {\mathrm{F}} \circ e _ {x} ^ {2} = (e _ {x} + i _ {x} ^ {\mathrm{F}}) ^ {2} \circ \lambda_ {\mathrm{F}} = (e _ {x} ^ {2} + \mathrm{F} (x, x)) \circ \lambda_ {\mathrm{F}},
\]

从而

\[
\begin{array}{r l} \lambda_ {\mathbf {F}} (x \otimes x \otimes u - \mathrm{Q} ^ {\prime} (x) u) & = (e _ {x} ^ {2} + \mathrm{F} (x, x) - \mathrm{Q} ^ {\prime} (x)) \circ \lambda_ {\mathbf {F}} (u) \\ & = (x \otimes x - \mathrm{Q} (x)) \otimes \lambda_ {\mathbf {F}} (u) \in \mathrm{I} (\mathrm{Q}). \end{array}
\]

引理 4. — 若二次型 Q 为零，则 C(Q) 不是别的，正是 E 的外代数。

  事实上，E 的外代数不是别的，正是 T(E) 对由形如 \(a \otimes v \otimes a\)（\(a \in E\)，\(v \in T(E)\)）的元素生成的双边理想 J 的商（第三章 § 5 第 5 小节与第 9 小节）。显然 I(Q) ⊂ J。因此只需证明 \(a \otimes v \otimes a \in I(Q)\)。若 \(v \in T^0\)，这是明显的。设这个断言对 \(n-1 \leqslant \nu \in \sum_{h=0}^{\infty} \mathrm{T}^{h}\) 已被证明，并设 \(x \in \mathrm{E}\) 与 \(u \in \mathrm{T}^{n-1}\)；我们有

\[
\begin{array}{c} \boldsymbol {a} \otimes \boldsymbol {x} \otimes \boldsymbol {u} \otimes \boldsymbol {a} = (\boldsymbol {a} + \boldsymbol {x}) \otimes (\boldsymbol {a} + \boldsymbol {x}) \otimes \boldsymbol {u} \otimes \boldsymbol {a} \\ - \boldsymbol {a} \otimes \boldsymbol {a} \otimes \boldsymbol {u} \otimes \boldsymbol {a} - \boldsymbol {x} \otimes \boldsymbol {a} \otimes \boldsymbol {u} \otimes \boldsymbol {a} - \boldsymbol {x} \otimes \boldsymbol {x} \otimes \boldsymbol {u} \otimes \boldsymbol {a} \end{array}
\]

  而右边的四项都属于 I(Q)。

  特别地，假设模 E 承认一个基 \((x_{i})_{i\in \mathbf{L}}\)，并让我们把指标集 L 全序化。我们知道（第三章 § 5 第 6 小节）E 的外代数承认由元素 \(x_{\mathrm{H}}\) 组成的族为基，其中 H 遍历 L 的有限子集所成之集，而 \(x_{\mathrm{H}}\) 是元素 \(x_{h_1}\wedge \ldots \wedge x_{h_q}\)，\((h_1,\ldots ,h_q)\) 表示 H 的元素的严格递增序列。

  另一方面，考虑由下列条件定义的双线性型 F：\(\mathrm{F}(x_{i}, x_{j}) = -\Phi(x_{i}, x_{j})\)（若 i \textgreater{} j），\(\mathrm{F}(x_{i}, x_{j}) = 0\)（若 i \textless{} j），而 \(\mathrm{F}(x_{i}, x_{i}) = -\mathrm{Q}(x_{i})\)。显然 \(\mathrm{Q}(x) + \mathrm{F}(x, x) = 0\)；用命题 3 的记号，由这个命题与引理 4 得出 \(\bar{\lambda}_{\mathrm{F}}\) 是 \(\wedge \mathrm{E}\) 到 C(Q) 上的一个同构，因而 C(Q) 是一个自由 A-模。让我们对 H 的元素个数作归纳来证明：我们有

\[
\overline {{{{\lambda}}}} _ {\mathrm{F}} (x _ {\mathrm{H}}) = \rho (x _ {h _ {1}}) \dots \rho (x _ {h _ {q}})\tag{11}
\]

\((\mathrm{where}\ (h_{1},\ldots,h_{q})\) 是 H 的元素的严格递增序列）。若 H 有 0 个或 1 个元素，这是明显的。设 (11) 对至多有 \(q-1\) 个元素的集合成立。考虑一个有 \(q\) 个元素的子集 H，用 \(j\) 记它的最小元素，并让我们写 \(\mathrm{H}=\{j\}\cup\mathrm{K}\)，其中 K 是有 \(q-1\) 个元素的子集。由 (8) 与归纳假设，我们有

\[
\overline {{{\lambda}}} _ {\mathrm{F}} (x _ {\mathrm{H}}) = \overline {{{\lambda}}} _ {\mathrm{F}} (x _ {j} \wedge x _ {\mathrm{K}}) = \rho (x _ {j}) \overline {{{\lambda}}} _ {\mathrm{F}} (x _ {\mathrm{K}}) + i _ {x _ {j}} ^ {\mathrm{F}} (\overline {{{\lambda}}} _ {\mathrm{F}} (x _ {\mathrm{K}})) = x _ {\mathrm{H}} ^ {\prime} + i _ {x _ {j}} ^ {\mathrm{F}} (x _ {\mathrm{K}} ^ {\prime}),
\]

这里对 L 的每个有限子集 J，置 \(x_{J}^{\prime}=\rho(x_{j_{1}})\ldots\rho(x_{j_{s}})\)，其中 \((j_{1},\ldots,j_{s})\) 是 J 的元素的递增序列。现在，对 \(i\in\mathbf{K}\)，\(x_{i}\) 属于线性型 \(y\to\mathrm{F}(x_{j},y)\) 的核，故 \(i_{x_{j}}^{\mathrm{F}}(x_{\mathrm{K}}^{\prime})=0\)（引理 1）。这就证明了所要的结果。

  这样我们就证明了下面的定理：

定理 1. --- 假设 A-模 E 承认一个基 \((x_{i})_{i\in L}\)，指标集 L 带有一个全序结构。对 L 的每个有限子集 H，置 \(x_{\mathrm{H}} = \rho(x_{h_1})\rho(x_{h_2})\ldots \rho(x_{h_q})\)，其中 \((h_1,\ldots,h_q)\) 是 H 的元素的严格递增序列。则诸元素 \(x_{\mathrm{H}}\) 构成 A-模 C(Q) 的一个基。

推论 1. --- 若 E 是维数为 n 的自由模，则 C(Q) 是维数为 \(2^{n}\) 的自由模；此外，若 \(n > 0\)，则 C \(^{+}\) 与 C \(^{-}\) 是维数为 \(2^{n-1}\) 的自由模。

  这由二项式系数的性质立即得出。

推论 2. --- 若 E 是自由模，则从 E 到 C(Q) 中的典范映射 \(\rho\) 与从 A 到 C(Q) 中的映射 \(a \to a\) .1 都是单射。

推论 3. --- 假设 E 是两个自由子模 \(E_{1}\) 与 \(E_{2}\) 的直和。设 \(Q_{i}\) 是 Q 在 \(E_{i}\) 上的限制，\(p_{i}\) 是从 C(Q \(_{i}\) ) 到 C(Q) 中的典范映射（i = 1, 2）。则

线性映射 \(p\)（从 \(\mathrm{C(Q_{1})\otimes C(Q_{2})}\) 到 \(\mathrm{C(Q)}\) 中）由双线性映射 \((a,b)\to p_{1}(a)p_{2}(b)\)（从 \(\mathrm{C(Q_{1})\times C(Q_{2})}\) 到 \(\mathrm{C(Q)}\) 中）导出，是一个双射。

  事实上，只需考虑由 \(E_{1}\) 的一个基与 \(E_{2}\) 的一个基取并所得到的 E 的基。

推论 4. --- 在推论 3 的假设与记号下，进一步假设 \(E_{1}\) 与 \(E_{2}\) 正交，并在 \(C(Q_{1})\otimes C(Q_{2})\) 上转移 \(C(Q)\) 的代数结构（借助双射 p）。若 \(a_{i}\) 与 \(b_{i}\) 是 \(C(Q_{i})\) 的偶元素或奇元素（i=1, 2），则我们有 \((a_{1}\otimes a_{2})(b_{1}\otimes b_{2})=\varepsilon(a_{1}b_{1})\otimes(a_{2}b_{2})\)，其中 \(\varepsilon=1\)，除非 \(a_{2}\) 与 \(b_{1}\) 都是奇的，此时 \(\varepsilon=-1\)。

  事实上，只需证明 \(p_2(a_2)p_1(b_1) = \varepsilon p_1(b_1)p_2(a_2)\)，为此可以假设 \(p_2(a_2)\)（相应地，\(p_1(b_1)\)）是元素的一个乘积 \(x_1 \ldots x_h\)（相应地，\(y_1 \ldots y_k\)），这些元素属于 \(\rho_{\mathbb{Q}}(\mathbf{E}_2)\)（相应地，\(\rho_{\mathbb{Q}}(\mathbf{E}_1)\)）。由于 \(\mathbf{E}_1\) 与 \(\mathbf{E}_2\) 正交，我们有

\[
x _ {i} y _ {j} + y _ {j} x _ {i} = \Phi (x _ {i}, y _ {j}) = 0,
\]

  从而

\[
x _ {1} \dots x _ {h} y _ {1} \dots y _ {k} = (- 1) ^ {h k} y _ {1} \dots y _ {k} x _ {1} \dots x _ {h}.
\]

  如果略去 \(\mathbf{E}_1\) 与 \(\mathbf{E}_2\) 是自由模这一假设，推论 3 与 4 的结论仍然成立（参见习题 1）。

